

```latex
\documentclass[11pt]{article}
\usepackage{amsmath,amssymb,amsfonts}
\usepackage{enumitem}
\usepackage{geometry}
\geometry{margin=1in}

\newcommand{\bk}{\Bbbk}
\newcommand{\cS}{\mathcal S}
\newcommand{\fb}{\mathfrak b}
\newcommand{\slq}{\mathfrak{sl}(n)}
\newcommand{\glq}{\mathfrak{gl}(n)}
\newcommand{\Uq}{U_q}
\newcommand{\bipr}{\bowtie}
\newcommand{\la}{\langle}
\newcommand{\ra}{\rangle}
\newcommand{\Kp}{K^+}
\newcommand{\Km}{K^-}
\newcommand{\id}{\mathrm{id}}
\newcommand{\End}{\mathrm{End}}
\newcommand{\op}{\mathrm{op}}

\title{The bigon algebra, the generalized quantum double of the Borel subalgebras, and \(U_q(\slq)\)}
\author{}
\date{}

\begin{document}

\maketitle

\begin{center}
\textbf{Convention.}
Throughout, \(\bk=\mathbb C(q^{1/n})\) (or any field of characteristic \(0\) containing \(q^{1/n}\)), with \(q\) transcendental, or else specialized to a value which is not a root of unity and satisfies \(q^2\neq 1\).  Put
\[
\lambda:=q^{-1}-q.
\]
All algebras are unital associative \(\bk\)-algebras.
\end{center}

\section{The extended quantum group and its Borel subalgebras}

We use the standard presentation of the simply-connected extended quantum group \(\Uq^{ext}(\slq)\) as in Klimyk--Schm\"udgen \cite{KS}.  It is generated by
\[
E_i,\ F_i\quad(1\le i\le n-1),\qquad \widehat K_i^{\pm1}\quad(1\le i\le n),
\]
with relations
\begin{align}
&\widehat K_i\widehat K_j=\widehat K_j\widehat K_i,\qquad 
\widehat K_1\widehat K_2\cdots \widehat K_n=1, \label{eq:Kprod}\\
&\widehat K_iE_j\widehat K_i^{-1}
=q^{\delta_{ij}-\delta_{i,j+1}}E_j,\qquad
\widehat K_iF_j\widehat K_i^{-1}
=q^{-\delta_{ij}+\delta_{i,j+1}}F_j, \label{eq:Kaction}\\
&E_iF_j-F_jE_i
=\delta_{ij}\,
\frac{\widehat K_i^{-1}\widehat K_{i+1}-\widehat K_i\widehat K_{i+1}^{-1}}
{q^{-1}-q}, \label{eq:EF}\\
&\sum_{r=0}^{1-a_{ij}}(-1)^r
\binom{1-a_{ij}}{r}_q
E_i^{1-a_{ij}-r}E_jE_i^r=0,\qquad i\neq j, \label{eq:SerreE}\\
&\sum_{r=0}^{1-a_{ij}}(-1)^r
\binom{1-a_{ij}}{r}_q
F_i^{1-a_{ij}-r}F_jF_i^r=0,\qquad i\neq j, \label{eq:SerreF}
\end{align}
where \(a_{ij}=2\delta_{ij}-\delta_{|i-j|,1}\) is the Cartan matrix of type \(A_{n-1}\), and
\[
[m]_q=\frac{q^m-q^{-m}}{q-q^{-1}},\qquad
[m]_q!=[1]_q\cdots[m]_q,\qquad
\binom{m}{r}_q=\frac{[m]_q!}{[r]_q![m-r]_q!}.
\]

The upper and lower Borel subalgebras are
\[
\Uq(\fb_+)=\langle E_i,\widehat K_i^{\pm1}\rangle,\qquad
\Uq(\fb_-)=\langle F_i,\widehat K_i^{\pm1}\rangle.
\]

The standard Drinfeld--Tanisaki skew Hopf pairing
\[
\sigma:\Uq(\fb_+)\times \Uq(\fb_-)\longrightarrow \bk
\]
is determined on generators by
\begin{equation}\label{eq:pairing}
\la \widehat K_i,\widehat K_j\ra=q^{\frac1n-\delta_{ij}},\qquad
\la E_i,F_j\ra=\frac{\delta_{ij}}{q^{-1}-q},\qquad
\la \widehat K_i,F_j\ra=\la E_j,\widehat K_i\ra=0,
\end{equation}
and is extended by the skew Hopf-pairing rules; see \cite[Ch.~6, 9]{KS}.

\section{A presentation of the generalized quantum double}

Let
\[
\mathcal D:=\Uq(\fb_-)\bipr \Uq(\fb_+)
\]
be the generalized quantum double with respect to the pairing \eqref{eq:pairing}.  We write
\[
K_i^+:=\text{the Cartan generator }\widehat K_i\text{ in the }\Uq(\fb_+)\text{ copy},
\]
\[
K_i^-:=\text{the Cartan generator }\widehat K_i\text{ in the }\Uq(\fb_-)\text{ copy}.
\]
We also write \(E_i\) for the positive Chevalley generators in the \(\Uq(\fb_+)\)-factor and \(F_i\) for the negative Chevalley generators in the \(\Uq(\fb_-)\)-factor.

The following presentation is obtained by unfolding the definition of the generalized quantum double; it is the form of the double used in \cite{KS} and is equivalent to Tanisaki's presentation \cite{Tanisaki}.

\begin{lemma}[Presentation of the double]\label{lem:double}
The algebra \(\mathcal D\) is generated by
\[
E_i,\ F_i\ (1\le i\le n-1),\qquad (K_i^+)^{\pm1},\ (K_i^-)^{\pm1}\ (1\le i\le n),
\]
subject to the following relations.
\begin{enumerate}[label=(D\arabic*)]
\item Cartan relations:
\[
K_i^\pm K_j^\pm=K_j^\pm K_i^\pm,\qquad
K_i^+K_j^-=K_j^-K_i^+,\qquad
\prod_{i=1}^n K_i^+=1,\qquad \prod_{i=1}^n K_i^-=1.
\]
\item Positive Borel relations:
\[
K_i^+E_j(K_i^+)^{-1}=q^{\delta_{ij}-\delta_{i,j+1}}E_j,
\]
and the quantum Serre relations \eqref{eq:SerreE} for the \(E_i\).
\item Negative Borel relations:
\[
K_i^-F_j(K_i^-)^{-1}=q^{-\delta_{ij}+\delta_{i,j+1}}F_j,
\]
and the quantum Serre relations \eqref{eq:SerreF} for the \(F_i\).
\item Mixed Cartan--root relations:
\[
K_i^-E_j(K_i^-)^{-1}=q^{\delta_{ij}-\delta_{i,j+1}}E_j,
\qquad
K_i^+F_j(K_i^+)^{-1}=q^{-\delta_{ij}+\delta_{i,j+1}}F_j.
\]
\item Mixed Chevalley relations:
\[
E_iF_j=F_jE_i\quad(i\neq j),
\]
and
\begin{equation}\label{eq:EFdouble}
E_iF_i-F_iE_i
=
\frac{(K_i^-)^{-1}K_{i+1}^- - K_i^+(K_{i+1}^+)^{-1}}
{q^{-1}-q}.
\end{equation}
\end{enumerate}
\end{lemma}

\begin{proof}[Sketch of proof]
The generalized quantum double of two Hopf algebras \(H_-\) and \(H_+\) with a skew Hopf pairing \(\sigma:H_+\times H_-\to\bk\) is the algebra generated by copies of \(H_-\) and \(H_+\) with cross relations
\begin{equation}\label{eq:cross}
yx=\sigma(x_{(1)},y_{(1)})\,\sigma(x_{(3)},y_{(3)})^{-1}\,x_{(2)}y_{(2)},
\qquad x\in H_+,\,y\in H_-;
\end{equation}
see for instance \cite[\S 6.3]{KS} or \cite{Majid}.  Applying \eqref{eq:cross} to the pairing \eqref{eq:pairing} and the standard coproducts of \(\Uq^{ext}(\slq)\) gives the mixed relations (D4) and (D5).  The Cartan part of \eqref{eq:pairing} gives the commutation of \(K_i^+\) with \(K_j^-\); the scalar \(q^{1/n}\) in \(\la K_i^+,K_j^-\ra\) is compatible with the determinant relations \(\prod K_i^\pm=1\).  A direct calculation for \(x=E_i\), \(y=F_j\) yields \eqref{eq:EFdouble}; the case \(i\neq j\) gives \(E_iF_j=F_jE_i\).  The remaining relations are exactly the defining relations of the two Borel subalgebras.  This is precisely the Tanisaki presentation of the generalized quantum double; see \cite{Tanisaki}, \cite[\S 9.2]{KS}.
\end{proof}

\begin{remark}
The quotient of \(\mathcal D\) by the relations \(K_i^+=K_i^-\) will be the usual \(\Uq^{ext}(\slq)\).  The point of the double is that before this quotient there are two independent Cartan \(n\)-tuples.
\end{remark}

\section{Root vectors}

Inside the positive Borel of \(\mathcal D\), define positive root vectors \(E_{ij}\) for \(1\le i<j\le n\) recursively by
\[
E_{i,i+1}=E_i,
\qquad
E_{ij}=E_iE_{i+1,j}-q^{-1}E_{i+1,j}E_i\quad(j>i+1).
\]
Inside the negative Borel, define negative root vectors \(F_{ij}\) for \(1\le j<i\le n\) recursively by
\[
F_{i+1,i}=F_i,
\qquad
F_{ij}=F_{i,j+1}F_j-qF_jF_{i,j+1}\quad(i>j+1).
\]
These are the standard root vectors attached to the reduced word
\((1,2,\dots,n-1)\); their basic properties are proved in \cite{Lu,KS}.  In particular, the elements
\[
\{E_{ij}\}_{i<j},\quad \{K_i^+\}_{i=1}^n
\]
generate the positive Borel, and the elements
\[
\{F_{ij}\}_{i>j},\quad \{K_i^-\}_{i=1}^n
\]
generate the negative Borel.

\section{The bigon algebra \(\cS_n(\mathfrak B^*)\)}

We recall the standard presentation of the algebra \(\cS_n(\mathfrak B^*)\) generated by the right-oriented horizontal arcs \(b_{ij}\) and \(c_{ij}\).  Equivalently, this is the FRT presentation of the quantum double of the two Borel subalgebras; see Faddeev--Reshetikhin--Takhtajan \cite{FRT}, Joseph \cite{Jo}, and Klimyk--Schm\"udgen \cite[\S 9]{KS}.  In the diagrammatic language, the generators \(b_{ij}\) are the lower horizontal arcs and the generators \(c_{ij}\) are the upper horizontal arcs.

Let \(V=\bk^n\) with basis \(e_1,\dots,e_n\), and let \(e_{ij}\in\End(V)\) be the matrix units.  Let \(\mathsf R\in\End(V\otimes V)\) be the normalized trigonometric \(R\)-matrix associated to the pairing \eqref{eq:pairing}; it is \(q^{1/n}\) times the inverse of the usual Jimbo \(R\)-matrix:
\begin{equation}\label{eq:Rmat}
\mathsf R
=
q^{1/n}\left(
q^{-1}\sum_i e_{ii}\otimes e_{ii}
+\sum_{i\neq j}e_{ii}\otimes e_{jj}
+(q^{-1}-q)\sum_{i<j}e_{ij}\otimes e_{ji}
\right).
\end{equation}
This normalization gives diagonal coefficients \(q^{1/n-\delta_{ij}}\), matching the Cartan part of \eqref{eq:pairing}.

For an \(n\times n\) matrix \(X=(x_{ij})\), write \(X_1=X\otimes I\) and \(X_2=I\otimes X\).

\begin{definition}\label{def:bigon}
The algebra \(\cS_n(\mathfrak B^*)\) is generated by entries of two matrices
\[
B=(b_{ij}),\qquad C=(c_{ij}),
\]
where
\[
b_{ij}=0\quad(i<j),\qquad c_{ij}=0\quad(i>j),
\]
the diagonal entries \(b_{ii},c_{ii}\) are invertible, and
\[
\prod_{i=1}^n b_{ii}=1,\qquad \prod_{i=1}^n c_{ii}=1,
\]
subject to the RTT relations
\begin{align}
\mathsf R C_1 C_2&=C_2C_1\mathsf R, \label{eq:RTC}\\
\mathsf R B_1 B_2&=B_2B_1\mathsf R, \label{eq:RTB}\\
\mathsf R C_1 B_2&=B_2C_1\mathsf R. \label{eq:RTCB}
\end{align}
\end{definition}

\begin{remark}
The relations \eqref{eq:RTC}--\eqref{eq:RTCB} are precisely the algebraic translation of the local relations for right-oriented horizontal arcs in the punctured bigon.  The triangularity conditions correspond to the fact that the \(b\)-arcs point to the lower side and the \(c\)-arcs to the upper side.  The determinant relations \(\prod b_{ii}=\prod c_{ii}=1\) are the \(\slq\) (rather than \(\glq\)) normalization.
\end{remark}

The following standard Gauss-decomposition theorem is the bridge between the FRT presentation and the Chevalley presentation.

\begin{lemma}[Gauss--FRT decomposition]\label{lem:Gauss}
Inside \(\cS_n(\mathfrak B^*)\), define
\[
K_i^+:=c_{ii},\qquad K_i^-:=b_{ii}^{-1},
\]
\[
E_i:=\lambda^{-1}c_{i,i+1}c_{i+1,i+1}^{-1},\qquad
F_i:=\lambda^{-1}b_{ii}^{-1}b_{i+1,i},
\qquad \lambda=q^{-1}-q.
\]
Then the elements \(K_i^\pm,E_i,F_i\) satisfy the relations (D1)--(D5) of Lemma \ref{lem:double}.

Conversely, in any algebra satisfying the relations (D1)--(D5), define matrices \(B=(b_{ij})\) and \(C=(c_{ij})\) by
\begin{equation}\label{eq:bcentries}
\begin{aligned}
c_{ii}&=K_i^+, & c_{ij}&=\lambda E_{ij}K_j^+\quad(i<j), & c_{ij}&=0\quad(i>j),\\
b_{ii}&=(K_i^-)^{-1}, & b_{ij}&=\lambda (K_j^-)^{-1}F_{ij}\quad(i>j), & b_{ij}&=0\quad(i<j).
\end{aligned}
\end{equation}
Then \(B\) and \(C\) satisfy the relations \eqref{eq:RTC}--\eqref{eq:RTCB} and the determinant relations.
\end{lemma}

\begin{proof}[Proof sketch]
This is the standard Gauss decomposition of the FRT algebra; see \cite{FRT,Jo,KS}.  Expanding the RTT relations in matrix coefficients gives:
\begin{itemize}
\item from \eqref{eq:RTC}, the positive Borel relations for \(K_i^+\) and the \(E_i\), including the quantum Serre relations;
\item from \eqref{eq:RTB}, the negative Borel relations for \(K_i^-\) and the \(F_i\);
\item from \eqref{eq:RTCB}, the mixed Cartan--root relations and the mixed Chevalley relation \eqref{eq:EFdouble}.
\end{itemize}
Conversely, substituting the expressions \eqref{eq:bcentries} into the RTT relations gives identities which follow from the defining relations (D1)--(D5) and the recursive definitions of the root vectors.  This is a finite matrix-coefficient calculation; it is written out in detail in \cite[\S 9.3--9.4]{KS} and in \cite{Jo}.
\end{proof}

\section{Explicit isomorphism \(\cS_n(\mathfrak B^*)\cong \Uq(\fb_-)\bipr\Uq(\fb_+)\)}

We can now state the main isomorphism.

\begin{theorem}\label{thm:main}
There is a unique algebra isomorphism
\[
\Phi:\cS_n(\mathfrak B^*)\xrightarrow{\sim}\mathcal D
=\Uq(\fb_-)\bipr\Uq(\fb_+)
\]
such that
\[
\Phi(c_{ii})=K_i^+,\qquad
\Phi(b_{ii})=(K_i^-)^{-1},
\]
and, for \(i<j\),
\[
\Phi(c_{ij})=\lambda E_{ij}K_j^+,
\]
while for \(i>j\),
\[
\Phi(b_{ij})=\lambda (K_j^-)^{-1}F_{ij}.
\]
Equivalently, the inverse isomorphism \(\Psi:\mathcal D\to\cS_n(\mathfrak B^*)\) is given by
\[
\Psi(K_i^+)=c_{ii},\qquad
\Psi(K_i^-)=b_{ii}^{-1},
\]
\[
\Psi(E_i)=\lambda^{-1}c_{i,i+1}c_{i+1,i+1}^{-1},\qquad
\Psi(F_i)=\lambda^{-1}b_{ii}^{-1}b_{i+1,i}.
\]
\end{theorem}

\begin{proof}
By Lemma \ref{lem:Gauss}, the assignments defining \(\Psi\) send the defining relations (D1)--(D5) of \(\mathcal D\) to relations holding in \(\cS_n(\mathfrak B^*)\).  Hence \(\Psi\) is a well-defined algebra homomorphism.

Conversely, again by Lemma \ref{lem:Gauss}, the matrices \(B\) and \(C\) defined inside \(\mathcal D\) by \eqref{eq:bcentries} satisfy the defining RTT relations of \(\cS_n(\mathfrak B^*)\).  Thus the assignments defining \(\Phi\) give a well-defined algebra homomorphism \(\Phi:\cS_n(\mathfrak B^*)\to\mathcal D\).

It remains to check that \(\Phi\) and \(\Psi\) are inverse.  On the Chevalley generators of \(\mathcal D\),
\[
\Phi\Psi(K_i^+)=\Phi(c_{ii})=K_i^+,\qquad
\Phi\Psi(K_i^-)=\Phi(b_{ii}^{-1})=K_i^-,
\]
\[
\Phi\Psi(E_i)
=
\lambda^{-1}\Phi(c_{i,i+1})\Phi(c_{i+1,i+1})^{-1}
=
\lambda^{-1}(\lambda E_iK_{i+1}^+)(K_{i+1}^+)^{-1}=E_i,
\]
\[
\Phi\Psi(F_i)
=
\lambda^{-1}\Phi(b_{ii}^{-1})\Phi(b_{i+1,i})
=
\lambda^{-1}K_i^-(\lambda (K_i^-)^{-1}F_i)=F_i.
\]
Since the \(E_i,F_i,K_i^\pm\) generate \(\mathcal D\), \(\Phi\Psi=\id_{\mathcal D}\).

Similarly, on the diagonal generators of \(\cS_n(\mathfrak B^*)\),
\[
\Psi\Phi(c_{ii})=\Psi(K_i^+)=c_{ii},\qquad
\Psi\Phi(b_{ii})=\Psi((K_i^-)^{-1})=b_{ii}.
\]
For the simple off-diagonal generators,
\[
\Psi\Phi(c_{i,i+1})
=
\Psi(\lambda E_iK_{i+1}^+)
=
\lambda\bigl(\lambda^{-1}c_{i,i+1}c_{i+1,i+1}^{-1}\bigr)c_{i+1,i+1}
=c_{i,i+1},
\]
\[
\Psi\Phi(b_{i+1,i})
=
\Psi(\lambda (K_i^-)^{-1}F_i)
=
\lambda b_{ii}\bigl(\lambda^{-1}b_{ii}^{-1}b_{i+1,i}\bigr)
=b_{i+1,i}.
\]
The higher off-diagonal generators are obtained recursively from the simple ones by the same \(q\)-commutator formulas which define the root vectors \(E_{ij}\) and \(F_{ij}\).  Hence \(\Psi\Phi\) fixes every \(b_{ij}\) and every \(c_{ij}\), so \(\Psi\Phi=\id_{\cS_n(\mathfrak B^*)}\).

Therefore \(\Phi\) is an algebra isomorphism.
\end{proof}

\section{The quotient by \(b_{ii}c_{ii}-1\) and \(U_q(\slq)\)}

Let \(I\) be the two-sided ideal of \(\cS_n(\mathfrak B^*)\) generated by
\[
b_{ii}c_{ii}-1,\qquad 1\le i\le n.
\]
Under the isomorphism \(\Phi\), these elements become
\[
\Phi(b_{ii}c_{ii}-1)=(K_i^-)^{-1}K_i^+-1.
\]
Let \(J\) be the two-sided ideal of \(\mathcal D\) generated by
\[
(K_i^-)^{-1}K_i^+-1,\qquad 1\le i\le n.
\]
Then
\[
\cS_n(\mathfrak B^*)/I\cong \mathcal D/J.
\]

We now identify \(\mathcal D/J\).

\begin{lemma}\label{lem:quotient}
There is a canonical algebra isomorphism
\[
\mathcal D/J\cong \Uq^{ext}(\slq).
\]
Consequently,
\[
\cS_n(\mathfrak B^*)/I\cong \Uq^{ext}(\slq).
\]
\end{lemma}

\begin{proof}
In \(\mathcal D/J\), the relations \((K_i^-)^{-1}K_i^+=1\) imply
\[
K_i^+=K_i^-\quad(1\le i\le n).
\]
Write \(H_i\) for their common image.  The relations (D1) become
\[
H_iH_j=H_jH_i,\qquad \prod_{i=1}^n H_i=1.
\]
The relations (D2) and (D3), together with the mixed Cartan--root relations (D4), become
\[
H_iE_jH_i^{-1}=q^{\delta_{ij}-\delta_{i,j+1}}E_j,
\qquad
H_iF_jH_i^{-1}=q^{-\delta_{ij}+\delta_{i,j+1}}F_j.
\]
The mixed Chevalley relation \eqref{eq:EFdouble} becomes
\[
E_iF_i-F_iE_i
=
\frac{H_i^{-1}H_{i+1}-H_iH_{i+1}^{-1}}{q^{-1}-q},
\]
and for \(i\neq j\),
\[
E_iF_j=F_jE_i.
\]
Together with the Serre relations for the \(E_i\) and \(F_i\), these are exactly the defining relations \eqref{eq:Kprod}--\eqref{eq:SerreF} of \(\Uq^{ext}(\slq)\).

Thus there is a surjective homomorphism
\[
\mathcal D/J\longrightarrow \Uq^{ext}(\slq)
\]
sending \(H_i\mapsto \widehat K_i\), \(E_i\mapsto E_i\), \(F_i\mapsto F_i\).  Conversely, the universal property of \(\Uq^{ext}(\slq)\) gives a homomorphism
\[
\Uq^{ext}(\slq)\longrightarrow \mathcal D/J
\]
sending \(\widehat K_i\mapsto H_i\), \(E_i\mapsto E_i\), \(F_i\mapsto F_i\).  These two homomorphisms are inverse on generators, hence are inverse isomorphisms.
\end{proof}

Combining Theorem \ref{thm:main} and Lemma \ref{lem:quotient} gives the desired result.

\begin{theorem}
There is a canonical algebra isomorphism
\[
\Uq(\slq)\cong \cS_n(\mathfrak B^*)/I,
\]
where \(I\) is the two-sided ideal generated by
\[
b_{ii}c_{ii}-1,\qquad 1\le i\le n.
\]
\end{theorem}

\begin{proof}
By Theorem \ref{thm:main},
\[
\cS_n(\mathfrak B^*)\cong \Uq(\fb_-)\bipr\Uq(\fb_+).
\]
Under this isomorphism, \(I\) corresponds to the ideal \(J\) generated by \((K_i^-)^{-1}K_i^+-1\).  By Lemma \ref{lem:quotient},
\[
(\Uq(\fb_-)\bipr\Uq(\fb_+))/J\cong \Uq^{ext}(\slq).
\]
Since \(\Uq^{ext}(\slq)\) is the simply-connected form of \(U_q(\slq)\) used in \cite{KS}, this proves the claim.
\end{proof}

\begin{thebibliography}{99}
\bibitem{Dr}
V.~G.~Drinfeld, \emph{Quantum groups}, Proceedings of the International Congress of Mathematicians, Berkeley, 1986, Amer. Math. Soc., 1987.

\bibitem{FRT}
L.~D.~Faddeev, N.~Yu.~Reshetikhin, L.~A.~Takhtajan, \emph{Quantization of Lie groups and Lie algebras}, Algebra i Analiz \textbf{1} (1989), 178--206; English transl. in Leningrad Math. J. \textbf{1} (1990), 193--225.

\bibitem{Jo}
A.~Joseph, \emph{Quantum Groups and Their Primitive Ideals}, Springer-Verlag, 1995.

\bibitem{KS}
A.~Klimyk, K.~Schm\"udgen, \emph{Quantum Groups and Their Representations}, Kluwer Academic Publishers, 1997.

\bibitem{Lu}
G.~Lusztig, \emph{Introduction to Quantum Groups}, Progress in Mathematics, vol.~110, Birkh\"auser, 1993.

\bibitem{Majid}
S.~Majid, \emph{Foundations of Quantum Group Theory}, Cambridge University Press, 1995.

\bibitem{Tanisaki}
T.~Tanisaki, \emph{Defining relations of quantized enveloping algebras}, in: Quantum Groups and Their Representations, Kluwer, 1997.
\end{thebibliography}

\end{document}
```